\section*{第 \ref{ch:g9:statproba} 章 --- 统计与概率}

\begin{solution}{exo:g9:statproba:1}
平均数：$\dfrac{14 + 9 + 17 + 9 + 12 + 11 + 12}{7} = \dfrac{84}{7} = 12$。

中位数：排序后为 $9, 9, 11, 12, 12, 14, 17$；中间（第 $4$ 个）数值为
$12$。

极差：$17 - 9 = 8$。
\end{solution}

\begin{solution}{exo:g9:statproba:2}
平均数：$\dfrac{19 + 21 + 24 + 18 + 22 + 25 + 19}{7} = \dfrac{148}{7}
\approx 21.1$ 度。

中位数：排序后为 $18, 19, 19, 21, 22, 24, 25$；第 $4$ 个数值为 $21$
度。
\end{solution}

\begin{solution}{exo:g9:statproba:3}
平均数：
\[
\frac{2 \times 0 + 4 \times 1 + 3 \times 2 + 1 \times 3}{10}
= \frac{0 + 4 + 6 + 3}{10} = \frac{13}{10} = 1.3
\text{ 每场进球数。}
\]
中位数：排序后的数据为 $0,0,1,1,1,1,2,2,2,3$；第 $5$ 个与第 $6$ 个
数值都是 $1$，中位数为 $1$。
\end{solution}

\begin{solution}{exo:g9:statproba:4}
$\P(3) = \frac16$. \quad
$\P(\text{偶数}) = \frac36 = \frac12$。 \quad
$\P(\text{至少 } 3) = \P(\{3,4,5,6\}) = \frac46 = \frac23$。 \quad
$\P(\text{不是 } 6) = 1 - \frac16 = \frac56$。
\end{solution}

\begin{solution}{exo:g9:statproba:5}
$\P(\text{白色}) = \frac{5}{12}$。

$\P(\text{红色或绿色}) = \frac{4 + 3}{12} = \frac{7}{12}$。

$\P(\text{不是白色}) = 1 - \frac{5}{12} = \frac{7}{12}$。

最后两个结果相等：“不是白色”\emph{就是}事件“红色或
绿色”。
\end{solution}

\begin{solution}{exo:g9:statproba:6}
这 $12$ 种结果（硬币的面，骰子的点数）
等可能。

正面朝上且掷得 $6$：$12$ 种结果中只有一种，即 $\frac{1}{12}$。

正面朝上或掷得 $6$：正面朝上的结果（共 $6$ 种）加上
（反面，$6$）这一结果，共 $7$ 种，故概率为 $\frac{7}{12}$。另一种方法是分别计算两事件，
再减去重叠部分，结果相同：
$\frac{6}{12} + \frac{2}{12} - \frac{1}{12} = \frac{7}{12}$.
\end{solution}

\begin{solution}{exo:g9:statproba:7}
第一次抽取：红球 $\frac25$，黑球 $\frac35$。第二次不放回抽取，
若第一次抽到红球，则剩下 $1$ 个红球、$3$ 个黑球，共 $4$ 个；
若第一次抽到黑球，则剩下 $2$ 个红球、$2$ 个黑球，共 $4$ 个。

两次都是黑球：$\frac35 \times \frac24 = \frac{6}{20} = \frac{3}{10}$。

颜色不同（红黑或黑红）：
$\frac25 \times \frac34 + \frac35 \times \frac24
= \frac{6}{20} + \frac{6}{20} = \frac{12}{20} = \frac35$.
\end{solution}

\begin{solution}{exo:g9:statproba:8}
\emph{1.} $9 \times 14 = 126$ 分。

\emph{2.} 平均每场 $15$ 分，$10$ 场比赛的总分为 $150$ 分，
她还需得 $150 - 126 = 24$ 分。
\end{solution}

\begin{solution}{exo:g9:statproba:9}
\emph{1.} 在 $36$ 个等可能的数对中，点数相同的数对为
$(1,1), (2,2), \dots, (6,6)$：共六个，所以
$\P(\text{点数相同}) = \frac{6}{36} = \frac16$。

\emph{2.} 两局游戏相互独立，树状图的每一层都分为赢（$\frac16$）与
输（$\frac56$）两支。从未赢过的概率为：
$\frac56 \times \frac56 = \frac{25}{36}$。于是
\[
\P(\text{至少赢一次}) = 1 - \frac{25}{36} = \frac{11}{36}
\approx 0.31 .
\]
\end{solution}

\begin{solution}{pb:g9:statproba:1}
\textbf{1.} 每次掷骰子，非六点的结果有 $5$ 种，共 $6$ 种；
$6 \times 6$ 表格显示，不含六点的结果有 $5 \times 5 = 25$ 种，
占全部 $36$ 种：$P(\text{无六点}) = \frac{25}{36}$，所以
$P(\text{至少一次}) = 1 - \frac{25}{36} = \frac{11}{36}
\approx 0.31$。

\textbf{2.} $\frac13 = \frac{12}{36} \neq \frac{11}{36}$。此处的加法
把“先六后六”这一结果在每个
$\frac16$ 中都计入一次，即多算了一次。可以
\emph{同时}发生的事件，其概率不能简单相加。

\textbf{3.} 和为 $7$ 的六种结果为 $(1,6), (2,5), (3,4),
(4,3), (5,2), (6,1)$：$P = \frac{6}{36} = \frac16$。和为 $2$ 的
只有 $(1,1)$：$\frac{1}{36}$。和为七的
出现方式最多 --- 正是和表格中最“厚”的中间部分。

\textbf{4.} 和为 $9$：$(3,6), (4,5), (5,4), (6,3)$ --- 四
种结果，$\frac{4}{36}$。和为 $10$：$(4,6), (5,5), (6,4)$ ---
三种，$\frac{3}{36}$。和为九的可能性更大：起决定作用的是计数，
而非用不计次序的面所写出和的“写法”数目，
即无序面组合的个数。

\textbf{5.} 树形图每一层都把不出六点的概率
乘以 $\frac56$，无论之前发生了什么：掷 $n$
次后，$P(\text{一次六点也未出现}) = \left(\frac56\right)^n$，
由对立事件法则即得该公式。

\textbf{6.} $\left(\frac56\right)^4 = \frac{625}{1296} \approx
0.482$，因此骑士获胜的概率为
 $1 - 0.482 = 0.518$ --- 每局平稳而持久的优势，接近
 $2\,\%$：一笔极好的赌注，而且它为他
付清了马车钱。

\textbf{7.} 掷 $7$ 次时，他的规则给出 $\frac76$ --- 一个
大于 $1$ 的概率，荒谬至极。这正是第 2 题
重复计数被放大：各次掷骰的成功彼此重叠，而
把它们的概率相加，就会把重叠部分一再计入。

\textbf{8.} $1 - \left(\frac{35}{36}\right)^{24} \approx
1 - 0.5086 = 0.491$：低于一半。以等额赌注押注于一个
 $49.1\,\%$ 的事件，骑士约输 $2$ 局，
（每 $100$ 局平均）--- 慢慢地、神秘地（对他而言），并且
必然地。

\textbf{9.} $1 - \left(\frac{35}{36}\right)^{25} \approx
1 - 0.494 = 0.506$：掷 $25$ 次时，这一赌注转为有利。
仅差一次掷骰，就决定了骑士能否盈利。

\textbf{10.} 把一个概率乘以尝试次数，
会把重叠的“成功”多次计入，最终
得出大于 $1$ 的荒谬答案。正确的做法
总是借助对立事件：将各次失败的概率
相乘，再用 $1$ 减去。

\textbf{11.} 平均数：$\frac{9 \times 2\,000 + 20\,000}{10} =
\frac{38\,000}{10} = 3\,800$ 欧元。中位数：中间两份工资
均为 $2\,000$，中位数为 $2\,000$ 欧元。这则广告
应当引用中位数 --- 十名工人中有九人从未见过
接近 $3\,800$ 的收入；平均数被一个
离群值拉高（见\cref{ex:g9:statproba:weighted}，已提醒注意权重）。

\textbf{12.} 例如 $2, 8, 15, 17, 18$：中位数为 $15$
（中间值），平均数为 $\frac{60}{5} = 12$。平均数低于
中位数，说明存在一段由\emph{低分}组成的尾部，把平均数
拉低，而上半部分的分数则偏高。

\textbf{13.} 前三个因子：
$1 \times \frac{364}{365} \times \frac{363}{365} \approx
0.992$。每加入一名新学生，就须多避开一个已被占用的日期，因此
各因子逐渐变小：$\frac{343}{365} \approx 0.94$ 对应
第 $23$ 位同学。整个乘积降至约 $0.493$ --- 于是
\[
P(\text{至少有两人生日相同}) \approx 0.507 :
\]
超过一半——在仅有 $23$ 人的班级里就是如此。直觉说
“$23$ 个人，$365$ 天，毫无机会”；而数学说
“抛一枚硬币”。

\textbf{14.} $\frac{23 \times 22}{2} = 253$ 对。每一对都是
一次新的巧合机会，而 $253$ 次机会
每次概率为 $\frac{1}{365}$，已不是小事：
一旦数的是对数而非人数，这种意外便消失了。

\textbf{15.} 示例方案：以四次为一组掷骰子，
记录每组是否出现六点，共 $100$ 组；或者
收集许多约 $23$ 人班级的生日名单，并
记录其中出现巧合的比例。观察到的
频率会有波动，但随着重复次数增加，
它应当越来越接近计算得到的概率
（$0.518$；$0.507$）--- 频率收敛于概率：
即大数定律，此处以期望的形式陈述，并
在高中分册中给出证明。
\end{solution}
